\documentclass[10pt,a4paper]{amsart}
\usepackage[margin=19mm]{geometry}
\usepackage{amsmath,amssymb,mathtools}
\usepackage[scr=rsfs,cal=euler]{mathalfa}
\usepackage{xfrac}
\usepackage[unicode,hidelinks,bookmarks=false]{hyperref}
\newcommand{\ie}{i.e.\ }
\newcommand{\cf}{cf.\ }
\newcommand{\resp}{respectively\ }
\newcommand{\Subsection}{\subsection}
\providecommand{\nobreakdash}{\nobreak\hbox{-}}
\setlength{\emergencystretch}{2em}
\multlinegap0pt
\usepackage{booktabs}
\usepackage{multirow}
\usepackage{amssymb}
\usepackage{mathtools}
\newtheorem{example}{Example}
\newtheorem{theorem}{Theorem}
\title{\vspace{-12mm} Subgradient Gliding Method for Nonsmooth Convex Optimization}
\author{Zhihan Zhu, Yanhao Zhang, Yong Xia}
\date{Source: arXiv:2602.14139v1 (CC BY 4.0)}
\begin{document}
\maketitle

\noindent\emph{Source and licence.} \vspace{-12mm} Subgradient Gliding Method for Nonsmooth Convex Optimization. Version arXiv:2602.14139v1 (CC BY 4.0), distributed by arXiv under the Creative Commons Attribution 4.0 International licence, http://creativecommons.org/licenses/by/4.0/, as recorded in the arXiv licence metadata for this exact version and shown on the abstract page https://arxiv.org/abs/2602.14139v1. Full licence notice: \texttt{licenses/arxiv-2602.14139-CC-BY-4.0.txt}.\par\smallskip

\noindent\textit{Editorial scope.} Counted unit: the article's theorem thm2.1 (source0.tex lines 334--336). The article's complete original proof of it is delivered with it (source0.tex lines 337--423, one \texttt{proof} environment of 87 lines). No mathematics is written, repaired or re-derived: the statement and the proof are the article's own lines, in the article's own notation, and no retained line is substituted or re-flowed.  Retained uncounted as the source-local context the proof needs: lines 251--253; lines 260--262; lines 294--296.  One declared adaptation: exactly 1 ancillary strings inside the retained spans are omitted (image-inclusion commands and, where the source repeats a label, the repeated label definitions) and no other line differs from the source; the figure environments, with their placement, captions and labels, are kept, and the package records the omitted strings in fidelity receipts bound to the affected bodies.  No retained line contains a citation, so the wrapper carries no bibliography and no bibliography entry.  The retained lines contain the article's own diagram code, which the wrapper typesets with the standard tikz packages; no image file is delivered and the package declares no asset dependency.  Editorial formatting only, in addition: an AMS \texttt{amsart} preamble that restates the article's own declarations and macros used by the retained lines, namely the environments \texttt{example}, \texttt{theorem} on separate counters, together with the standard packages those lines need.  The retained environments print in this document as Example 1, Theorem 1 in order of appearance, whereas the article prints the same items with the numbering of its own full document.  The complete native source of the article as distributed in the arXiv e-print for this exact version is bundled unchanged as \texttt{sources/194689/source0.tex}.
\begin{equation}
	\alpha_s=\frac{R}{\|g_s\|\sqrt{s}},~s=1,\cdots,t, \label{Nesterov}
\end{equation}

\begin{equation}
	\alpha_s=\frac{R}{G_s s^{\frac{a}{2}}}, ~s=1,\cdots,t,\label{size3}
\end{equation}

\begin{example}\label{e1}
	The function $f(x)=-\sqrt{r - k_1x(1)^2-k_2x(2)^2}$, where $r, k_1, k_2 > 0$, is a convex function on the ellipsoid $k_1x(1)^2+k_2x(2)^2\le r$, which lacks subgradients at boundary $k_1x(1)^2+k_2x(2)^2= r$.
\end{example}

\begin{theorem}[One-step failure of the classical PSG in convex setting]\label{thm2.1}
	In the worst case, when using the classical step sizes \eqref{Nesterov} and \eqref{size3} for convex objectives, PSG with the initial point drawn at random from $\text{int}(\mathcal{X})$ fails after a single iteration with probability arbitrarily close to 1.
\end{theorem}

\begin{proof}
	The construction of worst case in convex setting is based on Example \ref{e1}. We demonstrate that, as $k_2/k_1 \to \infty ~\text{or}~ 0$, the classical projected subgradient method with step sizes \eqref{Nesterov} and \eqref{size3} terminates after a single iteration with probability arbitrarily close to 1.
	
	Here, without loss of generality, we restrict attention to the case $k_2 > k_1$ and define $\rho := k_2 / k_1$. Since the optimal solution of the problem in Example \ref{e1} is $(0,0)$, it follows from the definition of $R$ that, for the step sizes \eqref{Nesterov} and \eqref{size3}, $R$ is chosen as the length of the major axis of the ellipse, namely $R = \sqrt{r/k_1}$. Moreover, in Example \ref{e1}, the objective function admits no subgradient at any point on the boundary $k_1 x_1^2 + k_2 x_2^2 = r$. Therefore, it suffices to show that, for sufficiently large $\rho$, the one-step iterate of the PSG with the classical step sizes \eqref{Nesterov} and \eqref{size3} lies outside the feasible set with arbitrary high probability.
	
	This is equivalent to showing that, for $x_1 \in \text{int}(\mathcal{X}) := \{ x : k_1 x_1^2 + k_2 x_2^2 < r \}$ and any $g_1 \in \partial f(x_1)$, the point $x_1 - \alpha_1 g_1$ does not belong to $\operatorname{int}(X)$. Noticing that the first step size of the classical rules \eqref{Nesterov} and \eqref{size3} coincides, namely, $\alpha_1 = R / \|g_1\|$, and
	\begin{equation*}
		g_1 =
		\begin{pmatrix}
			\frac{k_1 x_1(1)}{\sqrt{r - k_1 x_1(1)^{2} - k_2 x_1(2)^{2}}} \\
			\frac{k_2 x_1(2)}{\sqrt{r - k_1 x_1(1)^{2} - k_2 x_1(2)^{2}}}
		\end{pmatrix},
	\end{equation*}
	we have 
	\begin{eqnarray}
		x_1 - \alpha_1 g_1 &=&
		\begin{pmatrix}
			x_1(1) - R\frac{k_1 x_1(1)}{\sqrt{k_1^{2} x_1(1)^{2} + k_2^{2} x_1(2)^{2}}} \\
			x_1(2) - R\frac{k_2 x_1(2)}{\sqrt{k_1^{2} x_1(1)^{2} + k_2^{2} x_1(2)^{2}}}
		\end{pmatrix}\nonumber\\
		&=&
		\begin{pmatrix}
			x_1(1) - \sqrt{\frac{r}{k1}}\frac{x_1(1)}{\sqrt{x_1(1)^{2} + \rho^{2} x_1(2)^{2}}} \\
			x_1(2) - \sqrt{\frac{r}{k1}}\frac{\rho x_1(2)}{\sqrt{x_1(1)^{2} + \rho^{2} x_1(2)^{2}}}
		\end{pmatrix}. \nonumber~~~~({\rm by ~the~ definition~ of~  R~ and ~ \rho})
	\end{eqnarray}
	Hence, it suffices to show that
	\begin{equation*}
		k_1 \left(x_1(1) - \frac{\sqrt{\frac{r}{k1}}x_1(1)}{\sqrt{x_1(1)^{2} + \rho^{2} x_1(2)^{2}}}\right)^2 + \rho k_1 \left(x_1(2) - \frac{\sqrt{\frac{r}{k1}}\rho x_1(2)}{\sqrt{x_1(1)^{2} + \rho^{2} x_1(2)^{2}}}\right)^2 \geq r.
	\end{equation*}
	Since 
	\begin{eqnarray}
		&&k_1 \left(x_1(1) - \frac{\sqrt{\frac{r}{k1}}x_1(1)}{\sqrt{x_1(1)^{2} + \rho^{2} x_1(2)^{2}}}\right)^2 + \rho k_1 \left(x_1(2) - \frac{\sqrt{\frac{r}{k1}}\rho x_1(2)}{\sqrt{x_1(1)^{2} + \rho^{2} x_1(2)^{2}}}\right)^2 \nonumber\\
		&=& k_1 \left(x_1(1)^2 + \frac{r}{k1}\frac{x_1(1)^2}{x_1(1)^{2} + \rho^{2} x_1(2)^{2}} - 2\sqrt{\frac{r}{k1}}\frac{x_1(1)^2}{\sqrt{x_1(1)^{2} + \rho^{2} x_1(2)^{2}}}\right) + \nonumber\\
		&& \rho k_1 \left(x_1(2)^2 + \frac{r}{k1}\frac{\rho^2x_1(2)^2}{x_1(1)^{2} + \rho^{2} x_1(2)^{2}} - 2\sqrt{\frac{r}{k1}}\frac{\rho x_1(2)^2}{\sqrt{x_1(1)^{2} + \rho^{2} x_1(2)^{2}}}\right) \nonumber\\
		&=& k_1 x_1(1)^2 + \rho k_1 x_1(2)^2 + \frac{rx_1(1)^{2} + r \rho^3x_1(2)^{2}}{x_1(1)^{2} + \rho^{2} x_1(2)^{2}} - 2\sqrt{rk_1\left(x_1(1)^{2} + \rho^{2} x_1(2)^{2}\right)}, \nonumber
	\end{eqnarray}
	it suffices to show that, when $\rho$ is sufficiently large, for any $x_1 \in \text{int}(\mathcal{X}) := \{ x : k_1 x(1)^2 + k_2 x(2)^2 < r \}$, the probability that
	\begin{equation}\label{eq3}
		k_1 x_1(1)^2 + \rho k_1 x_1(2)^2 + \dfrac{r x_1(1)^2 + r \rho^3 x_1(2)^2}{x_1(1)^2 + \rho^2 x_1(2)^2} - 2 \sqrt{ r k_1 \left( x_1(1)^2 + \rho^2 x_1(2)^2 \right) } \ge r
	\end{equation}
	is arbitrarily close to 1.
	
	For any  $x_1 \in \text{int}(\mathcal{X}) := \{ x : k_1 x(1)^2 + k_2 x(2)^2 < r \}$ and $x_1 \neq 0$, we parameterize points in the elliptical region as follows:
	\begin{equation*}
		\begin{cases}
			x_1(1) = \sqrt{\frac{c}{k_1}} \cos \theta, \\
			x_1(2) = \sqrt{\frac{c}{\rho k_1}} \sin \theta,
		\end{cases}
	\end{equation*}
	where $c \in (0, ~r)$ and $\theta \in [0, ~2\pi)$. Substituting the parameterization into the left hand side of \eqref{eq3}, we obtain
	\begin{eqnarray}
		&&k_1 x_1(1)^2 + \rho k_1 x_1(2)^2 + \dfrac{r x_1(1)^2 + r \rho^3 x_1(2)^2}{x_1(1)^2 + \rho^2 x_1(2)^2} - 2 \sqrt{ r k_1 \left( x_1(1)^2 + \rho^2 x_1(2)^2 \right) } \nonumber\\
		&=& c + \frac{r\cos^2(\theta) + r\rho^2\sin^2(\theta)}{\cos^2(\theta)+\rho\sin^2(\theta)} - 2\sqrt{rc\left(\cos^2(\theta)+\rho\sin^2(\theta)\right)} \nonumber\\
		&=& c + \frac{r + r(\rho^2-1)\sin^2(\theta)}{1+(\rho-1)\sin^2(\theta)} - 2\sqrt{rc\left(1+(\rho-1)\sin^2(\theta)\right)} \nonumber\\
		&=& c + r + \frac{r\rho(\rho-1)\sin^2(\theta)}{1+(\rho-1)\sin^2(\theta)} - 2\sqrt{rc\left(1+(\rho-1)\sin^2(\theta)\right)}. \nonumber
	\end{eqnarray}
	Hence, \eqref{eq3} is equivalent to
	\begin{equation}\label{eq4}
		c + \frac{r\rho(\rho-1)\sin^2(\theta)}{1+(\rho-1)\sin^2(\theta)} - 2\sqrt{rc\left(1+(\rho-1)\sin^2(\theta)\right)} \ge 0.
	\end{equation}
	For any arbitrarily small $\epsilon > 0$, as long as $\theta \in [\epsilon, \pi-\epsilon] \cup [\pi + \epsilon, 2\pi-\epsilon]$, there exists a positive constant $M$ such that \eqref{eq4} holds for all $\rho > M$. This is because
	\begin{eqnarray}
		&&c + \frac{r\rho(\rho-1)\sin^2(\theta)}{1+(\rho-1)\sin^2(\theta)} - 2\sqrt{rc\left(1+(\rho-1)\sin^2(\theta)\right)} \nonumber\\
		&\geq& c + \frac{r\rho(\rho-1)\sin^2(\epsilon)}{1+(\rho-1)} - 2\sqrt{rc\left(1+(\rho-1)\right)} \nonumber\\
		&=& c + r(\rho-1)\sin^2(\epsilon) - 2\sqrt{rc\rho}. \label{eq5}
	\end{eqnarray}
	Since the second term in \eqref{eq5} is of linear order in $\rho$ while the third term is of order $1/2$ in $\rho$, \eqref{eq5} can be made arbitrarily large when $\rho$ is sufficiently large. Therefore, there must exist some positive number $M$ such that for all $\rho > M$, \eqref{eq5} is greater than zero, and consequently \eqref{eq4} holds. Figure \ref{fig7} illustrates this.
	\begin{figure}[h]
		\centering
		
		\vspace{2mm}
		\caption{Lower bound on the area of intial points that lead to one-step termination of PSG.}
		\label{fig7}
	\end{figure}

	This implies that when $\rho > M$, if the initial point $x_1$ is chosen from the region $\text{int}(\mathcal{X}) := \{ x : k_1 x(1)^2 + k_2 x(2)^2 < r \}$, and we denote by $\mathcal{X}_1$ the set of all $x_1$ that cause the PSG to terminate in one step, then
	\begin{equation*}
		\mathbb{P}\left(x_1 \in \mathcal{X}_1\right) = \dfrac{\text{area}(\mathcal{X}_1)}{\text{area}(\mathcal{X})} \geq \dfrac{\sqrt{\frac{r}{k_1}}\sqrt{\frac{r}{\rho k_1}}(\pi-2\epsilon)}{\sqrt{\frac{r}{k_1}}\sqrt{\frac{r}{\rho k_1}}\pi} = 1 - \frac{2\epsilon}{\pi}.
	\end{equation*}
%	\begin{eqnarray}
%		\mathbb{P}\left(x_1 \in \mathcal{X}_1\right) &=&  \frac{\text{area}(\mathcal{X}_1)}{\text{area}(\mathcal{X})} \nonumber \\
%		&\geq& \frac{\sqrt{\frac{r}{k_1}}\sqrt{\frac{r}{\rho k_1}}(\pi-2\epsilon)}{\sqrt{\frac{r}{k_1}}\sqrt{\frac{r}{\rho k_1}}\pi} \nonumber \\
%		&=& 1 - \frac{2\epsilon}{\pi}. \nonumber
%	\end{eqnarray}
	The theorem is thus proved by the arbitrariness of $\epsilon$.
\end{proof}

\end{document}
